\begin{center}\bfseries Résumé\end{center}
\begingroup
\fontsize{9.3}{10.4}\selectfont
\setlength{\parskip}{1.5pt plus .5pt minus .5pt}
\noindent
On construit un opérateur positif de masse sur les fibres de particules
persistantes et on détermine sa spécialisation scalaire pour l’électron
central. Le résultat relie l’identité de la particule, la classification
de ses états et l’évaluation de sa masse : la grandeur provient d’une
structure qui conserve la route, l’orientation et la mémoire de sa
génération. Dans le secteur électronique, deux lecteurs bidirectionnels
coïncident sur une trajectoire explicite et fixent la même correction de
retour ; les autres collections conservent la portée opératorielle, spectrale ou
scalaire démontrée pour leurs propres fibres.


La trajectoire centrale parcourt exactement trois diagonales nonadiques et
vingt-sept cellules. Deux histoires ayant le même histogramme de valeurs absolues
conservent des mémoires réversibles différentes. La transformée finie du mot
de retour localise trois modes non nuls et démontre spectralement
$(\Omega,P_6)=(80,6)$. Sur les multisections, la loi relative des masses
fournit un critère opératoriel nécessaire et suffisant : un transport TPK
bijectif conserve la masse exactement lorsqu’il commute avec l’opérateur de masse,
ou, de manière équivalente, lorsqu’il conserve sa coordonnée logarithmique complète.


Le recensement régional parcourt \(104\,976\) paires de germes, produit \(468\)
émissions et retient \(243\) mots distincts. La bijection tritique par paires
situe ces mots dans le corps \(S_0=E_9^3\) de \(729\) positions. La
différence exacte de seuil détermine alors la couronne
\(243+27+1=271\) et la complétion
\(729+243+27+1=1000\), avec une normalisation compatible sous les projections
entre dimensions. La sélection régionale agit ensuite : ni la bijection ni
la complétion cubique ne définissent une espèce ou une valeur métrologique.


Le cadre proposé est l’holographie modulaire triadique, ci-après HMT, et
sa réalisation au moyen de la mécanique dimensionnelle. Son noyau réunit trois objets :
\emph{All Possible Paths}, ci-après APP, est une structure arithmétique
de chemins sur \((\mathbb Z/9\mathbb Z)^2\) : son graphe de Cayley additif,
à générateurs cardinaux, fixe l’adjacence, et sa réalisation cellulaire
est toroïdale. Les opérations et les évaluations de somme et de produit se
distinguent au sein de cette structure ;
TRIT est l’unité de signature orientée \(\{-1,0,+1\}\), qui distingue les
régimes locaux ; le \emph{Topological Phase Kernel}, ci-après TPK,
compose la dynamique de sélection, de transport et de mémoire. Son évolution
enrichie conserve feuillets, résidus, quotients, frontière et retenue, et
produit la structure discrète conjointe du continu. La composition de
\(108\) mises à jour détermine le registre dodécaphasique, et la transformée
orbitale reconstruit réversiblement \(K\), antécédents des coordonnées utilisées par le lecteur
de masse.


La particule se définit par la persistance d’extensions compatibles ;
la route enregistre son histoire observable. L’atlas distingue \(81\) cellules,
\(56\) configurations de route, \(13\) multisections et \(324\) conditions
orientées. On construit leurs quotients, les observables de saveur et de couleur,
la section centrale unique et les opérations de conjugaison, d’hélicité et
d’holonomie, y compris le ruban de Möbius avec sa fibre et sa boucle. La signature
intégrale et les moments des canaux déterminent le caractère positif ;
l’action, l’horloge et la coordinatisation des unités fournissent la section
absolue. Les rapports de masse conservent le caractère relatif et les
différences entre exposants de retour.


L’incidence longitudinale marquée fixe d’abord
\(N=5760\), \(r_N=5759/23040\) et
\(L_\alpha=\alpha_{\rm HMT}^{16}r_NL_\star\). De cette longueur sont induits
l’horloge, \(R_{108}=L_\alpha\), l’énergie \(Q_L\), la masse chronologique et
la cotransformation du facteur électronique. Sur chaque fibre massive, les
identités \(H\Lambda=\hbar c_{\rm int}I\) et
\(\mathcal R_g\Lambda=L_\alpha^2I\) déterminent de manière unique
\(G=c_{\rm int}^3L_\alpha^2/\hbar\) ; son expression circulaire est un
corollaire ultérieur. La déduction n’utilise pas \(G\) comme antécédent de
l’opérateur de masse et ne modifie pas les tableaux du spectre.


La composition se réalise sur les registres avant d’évaluer leurs signatures.
Le développement comprend les secteurs de masse nulle, la propagation orientée,
les mélanges, les résonances et les réalisations orbitales de Kepler--Sommerfeld. La
complétion graduée produit les règles d’occupation et l’exclusion de
Pauli sur des modes complets, avec leurs étiquettes de spin, de feuillet, de saveur et de
couleur. La comparaison métrologique réunit chaque évaluation avec sa route,
son opérateur et sa coordinatisation, et distingue valeurs centrales, intervalles, limites et
largeurs. Ainsi sont reliées structure, équation et valeur en conservant le
contenu spécifique de chaque résultat et de son identification physique.


\noindent\textbf{Mots-clés :} dérivation du spectre de masses ;
structure des particules ; holographie modulaire triadique ; mécanique
dimensionnelle ; persistance ; holonomie ; conjugaison ; comparaison métrologique.

\par\endgroup
